\section{Conclusion}
\label{sec:conclusion}

The Dex programming model gives an application Steps that never block,
state that is read from a snapshot and written at a commit, Channels,
Timers, and a decision at the end of each Step. We have shown that this is
enough to express a conventional concurrent language: every well-formed
program of \Lang{}, with its threads, mutexes, condition variables, timed
receives, cancellation, and exceptions, translates into a Flow with the
same observable finite traces and the same fair maximal traces, in a
failure-free environment. We have also shown that, in the same fragment,
failures are transparent: crashes, takeovers, Worker retries, zombie
Workers, and duplicate deliveries add nothing observable except repeated
external effects, each equal to an earlier one. The statements are
precise, their scope is explicit, and their proofs are checked by the Lean
kernel from three standard axioms.

The constructive content is the translation. Each live thread that has not
halted is carried by one active Step, whose input is the thread's state; a
halted thread has none. A mutex is a Channel with one token. A
condition variable is a waiter list in an Attribute and one Channel per
waiter. A cancellation is a message that competes with the operation's own
condition in an \kw{AnyOf}. A timed operation passes through an Enter Step
and a shared timed Step with one Timer.
None of these needs a lock held across a wait or an atomic
read-modify-write, and the proof shows why none is missed.

Three lessons go beyond Dex. First, writing an operational semantics
against an implementation, rule by rule with citations, is a review of the
implementation in its own right; it found two defects before any proof
existed. Second, an exact equivalence forces decisions that informal
descriptions leave open, about time, about ill-behaved programs, and about
what the environment may send, and it is better to make them explicit than
to weaken the theorem until they disappear. Third, fixed statements make
proofs parallel, but a fixed simulation invariant must be inductive, and
that can be tested cheaply before anyone tries to prove it. We lost time to
six gaps that a fuzz of perturbed reachable states finds in under a
minute.

The failure theorem adds a lesson of its own. Repeated effects are not
arbitrary: because every attempt of a method reuses one snapshot and runs
the same handler, the first occurrences of its effects come in order, even
when zombies overlap with newer attempts. That is what makes erasure
work, and it is also why an application needs more than idempotent
effects: a late repeat can still overwrite a newer write, so effects need
idempotency keys, commutativity, or a fence by attempt.

The limits are equally explicit. The Dex side is a model validated by
review, and liveness assumes fair scheduling and Workers that settle. The
failure theorem covers SYNC methods only, assumes that no method exhausts
its retries and that user code ignores its attempt, keeps the timer order
as an idealization, and has a backward liveness inclusion only for words
without StartFlow retries. Extending the admissible environment and the
failure model with RPC entry points, SubFlows, and Attribute locks, and
connecting the model to the server's deterministic simulation tests, are
the next directions.
